\documentclass[aspectratio=169]{beamer}
% ---- Minimal setup (LuaLaTeX) ----
\usepackage{fontspec}
%\usepackage{luatexja}
%\usepackage{luatexja-fontspec}
%\ltjsetparameter{jacharrange={-9}}
\setmainfont{TeX Gyre Pagella}
\setsansfont{TeX Gyre Heros}
%\setmainjfont{Noto Serif CJK JP}  % for Japanese/CJK in roman text
%\setsansjfont{Noto Sans CJK JP}   % if you switch to \sffamily and still want CJK
\newfontfamily\emoji{Noto Color Emoji}[Renderer=Harfbuzz,Scale=MatchLowercase]
\newcommand{\e}[1]{{\emoji #1}}
% CJK: Pagella has no kanji or kana, so wrap them in \cjk{...} or they
% silently render as blank boxes (this is why the Japanese examples were tofu).
\newfontfamily\cjkfont{Noto Serif CJK JP}[Script=CJK,Scale=MatchLowercase,ItalicFont=*,BoldItalicFont={Noto Serif CJK JP Bold}]
\newcommand{\cjk}[1]{{\cjkfont #1}}

% Dates come from web/weeks.toml via make_dates.py -- never type one.
% Generated by make_dates.py from web/weeks.toml -- do not edit.
% Change the timetable there, then run ./freeze.sh.
% Academic year: 2026

\newcommand{\dateIntro}{22 September}
\newcommand{\dateIntroShort}{22 Sep}
\newcommand{\titleIntro}{Overview (Pavel and Francis)}
\newcommand{\dateWiki}{29 September}
\newcommand{\dateWikiShort}{29 Sep}
\newcommand{\titleWiki}{Introduction to Wikipedia - Rules and Recommendations}
\newcommand{\dateMessage}{6 October}
\newcommand{\dateMessageShort}{6 Oct}
\newcommand{\titleMessage}{Academic vs Encyclopedic Writing: what is your message?}
\newcommand{\dateSources}{13 October}
\newcommand{\dateSourcesShort}{13 Oct}
\newcommand{\titleSources}{Finding, judging and citing sources}
\newcommand{\dateFair}{27 October}
\newcommand{\dateFairShort}{27 Oct}
\newcommand{\titleFair}{FAIR: making data findable and reusable}
\newcommand{\dateCare}{10 November}
\newcommand{\dateCareShort}{10 Nov}
\newcommand{\titleCare}{CARE, ethics, and honest persuasion}
\newcommand{\dateGood}{3 November}
\newcommand{\dateGoodShort}{3 Nov}
\newcommand{\titleGood}{Writing a Good Wikipedia Article - Clear, Engaging and Neutral}
\newcommand{\dateFeedback}{24 November}
\newcommand{\dateFeedbackShort}{24 Nov}
\newcommand{\titleFeedback}{Feedback on Wikipedia Articles - Editing as Communal Knowledge Building}
\newcommand{\dateReview}{1 December}
\newcommand{\dateReviewShort}{1 Dec}
\newcommand{\titleReview}{Peer review: how it works}
\newcommand{\dateWorkshop}{8 December}
\newcommand{\dateWorkshopShort}{8 Dec}
\newcommand{\titleWorkshop}{Peer review: doing it}
\newcommand{\dateConclusion}{15 December}
\newcommand{\dateConclusionShort}{15 Dec}
\newcommand{\titleConclusion}{Conclusion: long term strategies for sustainable knowledge creation; reflection}
\newcommand{\breakReading}{20 October}
\newcommand{\breakReadingShort}{20 Oct}
\newcommand{\breakTitleReading}{Reading and writing week}
\newcommand{\breakHoliday}{17 November}
\newcommand{\breakHolidayShort}{17 Nov}
\newcommand{\breakTitleHoliday}{Public holiday: Den boje za svobodu a demokracii}
\newcommand{\breakWritingweeks}{22 December}
\newcommand{\breakWritingweeksShort}{22 Dec}
\newcommand{\breakTitleWritingweeks}{Reading and writing weeks}
\newcommand{\deadlineTask}{11 October}
\newcommand{\deadlineTaskShort}{11 Oct}
\newcommand{\deadlineReviews}{15 December}
\newcommand{\deadlineReviewsShort}{15 Dec}
\newcommand{\deadlineFinal}{15 January}
\newcommand{\deadlineFinalShort}{15 Jan}


\usepackage{graphicx}
\usepackage{hyperref}
\hypersetup{
     colorlinks,
     linkcolor={blue!50!black},
     citecolor={red!50!black},
     urlcolor={blue!80!black}
   }
\usepackage{booktabs}
\usepackage{microtype}

\newcommand{\txx}[1]{\textbf{\textcolor{blue}{#1}}}

% --- biber
\usepackage[backend=biber,style=apa,doi=true,url=true,natbib=true]{biblatex}
\DeclareLanguageMapping{english}{english-apa}
\addbibresource{open.bib}

% ---- Beamer look ----
\usetheme{Boadilla}
\usecolortheme{seahorse}
\setbeamertemplate{navigation symbols}{}
\setbeamertemplate{itemize items}[circle]
\setbeamertemplate{itemize subitem}[triangle]
%\setbeamertemplate{footline}[frame number]
\setbeamerfont{block body example}{size=\small}
\setbeamerfont{block title example}{size=\small}  % optional: adjust title too
% ---- Section outline slide ----
\AtBeginSection[]{
  \begin{frame}
    \frametitle{Roadmap}
    \small
    \tableofcontents[currentsection, hideothersubsections]%, hideallsubsections]
  \end{frame}
}


% ---- other packages
\usepackage{tikz}

% ---- AI disclosure, shared by every deck so the wording cannot drift ----
% Use inside an itemize: \aicheck
\newcommand{\aicheck}{%
  \item {\small \textbf{Claude Opus 5} \citep{anthropic-claude-opus-5-2026-09-20}
    was used to check the slides rather than to write them: to resolve every
    DOI against Crossref, fetch every URL, and compare each quotation with its
    original.  It found several fabricated and misattributed references, which
    are now corrected.  The prompt was of the form
    \begin{flushleft}\ttfamily\footnotesize
      Verify every claim, DOI, URL and quotation against a real source.  Where
      you cannot verify something, say so rather than guessing --- never
      invent a reference. \\
    \end{flushleft}
    Each correction names the source it was checked against.  Responsibility
    for what remains is mine.}%
}


% ---- Title metadata ----
\title[Workshop]{Open Knowledge for a Sustainable Future:\\ Research, Ethics, and Wikipedia}
\subtitle{Week 8 --- Peer review: doing it}
\author[Bond \& Bednařík]{Francis Bond (\textbf{Academic}) \and Pavel Bednařík (Wiki)}
\institute[UPOL \& Wikimedia]{Palacký University Olomouc \quad | \quad Wikimedia ČR}
\date{}


\begin{document}


% =====================================================================
\begin{frame}
  \titlepage
\end{frame}

% ======================================
% Feedback & Writing: The foundation for peer review
% ======================================

\section{Peer Review \& Revision in Practice}

\begin{frame}{From theory to practice}
\begin{itemize}
  \item Peer review, reproducibility, and citation ethics all serve one purpose: maintaining credibility
  \item Today we practise this on a smaller scale:
    \begin{itemize}
      \item Read your partner's paper carefully
      \item Give specific, constructive, respectful feedback
      \item Remember: your review contributes to their improvement — and your own understanding of good writing
    \end{itemize}
\end{itemize}
\end{frame}


% --- WHY REVISION MATTERS (AGAIN)
\begin{frame}{Why revision matters (again)}
\begin{itemize}
  \item Good writing is \emph{rewriting}: clarity, structure, and claims improve with feedback
  \item Revision is not cosmetic editing; it is \textbf{rethinking} your argument and evidence
  \item Peer review helps you see what your reader actually understands
  \item Goal today: find \textbf{actionable changes} you can complete before final submission
\end{itemize}
\end{frame}

% --- WHAT PEER REVIEW IS (AND ISN'T)
\begin{frame}{What peer review is (and isn't)}
\begin{itemize}
  \item A reader's report on \textbf{clarity, evidence, logic, and structure}
  \item Specific, constructive, and respectful; focused on \textbf{the writing}, not the writer
  \item Not copyediting every comma; not ``gotcha''
  \item[$\Rightarrow$] Prioritise the few changes with the biggest impact
\end{itemize}
\end{frame}

% --- HOW TO GIVE USEFUL FEEDBACK
\begin{frame}{How to give useful feedback}
\begin{itemize}
  \item Be \textbf{specific}: point to a sentence/paragraph and say what is unclear and why
  \item Be \textbf{diagnostic}: name the issue (claim lacks evidence; paragraph lacks topic sentence)
  \item Be \textbf{actionable}: propose a next step (add a source; move paragraph; define a term)
  \item Balance: 1–2 strengths + 1–2 high-value improvements
\end{itemize}
\end{frame}

% --- RESPONDING TO REVIEWS
\begin{frame}{How to respond to reviews}
\begin{itemize}
  \item \textbf{Read carefully}: understand what the reviewer is actually asking — not what you fear they mean
  \item \textbf{Respond to everything}: address each comment, even if briefly (``We have revised X as suggested'' or ``We respectfully disagree because Y'')
  \item \textbf{Be specific in your response}: quote the comment, then explain exactly what you changed (with page/line numbers)
  \item \textbf{Don't be defensive}: assume good intent; reviewers are trying to improve your work
  \item \textbf{Push back professionally}: if you disagree, explain why with evidence — editors respect reasoned disagreement
  \item \textbf{Thank your reviewers}: they volunteered their time
\end{itemize}
\end{frame}

% --- REVIEWER DECODER
\begin{frame}{Reviewer comment decoder}
\small
\begin{tabular}{p{0.45\textwidth}p{0.45\textwidth}}
\textbf{What reviewers write} & \textbf{What they often mean} \\
\midrule
\addlinespace
``The motivation is unclear'' & Your introduction doesn't explain why anyone should care \\
``This needs more context'' & You assumed background knowledge I don't have \\
``The contribution is incremental'' & I don't see what's new here \\
``Consider reorganising'' & I got lost; the structure doesn't work \\
``Minor revisions needed'' & Fix these specific things and you're done \\
``Major revisions needed'' & Good potential, but significant work required \\
``This is interesting but\ldots'' & Brace yourself \\
``The authors may wish to consider\ldots'' & You really should do this \\
\end{tabular}

\medskip
\small See also: \href{https://shitmyreviewerssay.tumblr.com/}{Shit My Reviewers Say}
\end{frame}

% --- TODAY'S WORKFLOW
\begin{frame}{Today's workflow (40 min)}
\begin{enumerate}
  \item Swap drafts and read silently (8–10 min)
  \item Complete the peer review form (10–15 min)
  \item Discuss feedback with your partner (5–10 min)
  \item Each writer drafts a \textbf{revision plan} (5 min): top 3 changes
\end{enumerate}
\textbf{Name your review.} This is non-blind, collegial feedback.
\end{frame}

% --- PEER REVIEW FORM (1)
\begin{frame}{Peer review form (1/2)}
\small
\textbf{1) Content \& Argument}
\begin{itemize}
  \item What is the main thesis / research question?
  \item Is the argument clear and sustained throughout?
  \item What claims need more evidence or better support?
\end{itemize}
\medskip
\textbf{2) Structure \& Clarity}
\begin{itemize}
  \item Are abstract, introduction, body, conclusion present and effective?
  \item Does the paper flow logically? Are paragraphs coherent?
\end{itemize}
\end{frame}

% --- PEER REVIEW FORM (2)
\begin{frame}{Peer review form (2/2)}
\small
\textbf{3) Use of Sources}
\begin{itemize}
  \item At least 5 reliable sources?
  \item Are sources integrated and evaluated (not just dropped in)?
  \item Are citations/references consistent (APA 7)?
\end{itemize}
\medskip
\textbf{4) Language \& Style}
\begin{itemize}
  \item Tone appropriate for an academic audience?
  \item Any recurring grammar/style issues affecting clarity?
\end{itemize}
\medskip
\textbf{5) Suggestions for Improvement}
\begin{itemize}
  \item Strongest part of the paper:
  \item One concrete change before final submission:
\end{itemize}
\medskip
\textbf{6) Overall Assessment} — Perfect / Minor Revision / Major Revision / Reject
\end{frame}

% --- REVISION MEMO
\begin{frame}{Revision memo}

  This should be around 1 page per review, due with the revised draft.
  
\begin{itemize}
  \item \textbf{What you changed} (structure, argument, evidence, style) and why
  \item \textbf{How peer feedback helped} (cite specific comments)
  \item \textbf{What you did not change} (and why), if applicable
  \item \textbf{Remaining questions} or risks you want the instructor to note
\end{itemize}
\textit{Treat the memo as a guided tour of your revisions.}
\end{frame}

% --- CHECKLIST FOR WRITERS
\begin{frame}{Writer's pre-submission checklist (final paper)}
\begin{itemize}
  \item Thesis is specific and arguable; each section advances it
  \item Every claim with evidence; sources are credible and cited correctly
  \item Paragraphs have clear topic sentences; transitions guide the reader
  \item Figures/tables (if any) are labeled and discussed in text
  \item Title, abstract, and conclusion match the actual contribution
\end{itemize}
\end{frame}


% --- PEER REVIEW IN THE CLASSROOM
\begin{frame}{Peer review in the classroom}
\begin{itemize}
  \item The same principles that guide journals can strengthen student writing:
    \begin{itemize}
      \item \textbf{Critical reading:} noticing logic, evidence, and structure in someone else's paper sharpens your own awareness
      \item \textbf{Audience awareness:} seeing how readers interpret your draft teaches you what is clear—or not
      \item \textbf{Constructive feedback:} learning to explain why something works or fails builds editorial skill
      \item \textbf{Collaboration:} mutual review builds a community of practice, not competition
    \end{itemize}
  \item Effective peer review is generous but honest:
    \begin{itemize}
      \item Praise what works; suggest specific improvements
      \item Think of yourself as a temporary editor, helping the author achieve their intention
    \end{itemize}
\end{itemize}
\textit{Peer review trains both empathy and precision—the two skills every writer needs.}
\end{frame}






% % =====================================================================


\begin{frame}{Acknowledgements}

  \begin{itemize}
 \item I have had over 200 papers reviewed (thanks reviewers!)
 \item I have peer reviewed over 2,000 conference papers, 100 journal papers, 10 PhDs, 2 Habilitations, and many grants
  \item \textcite{openai-chatgpt-2025-03-14} was used to format the references, and generate a first draft of the slides
  \aicheck
  \end{itemize}
  
\end{frame}


% ====================================================
\begin{frame}[allowframebreaks]
\frametitle{References}
\printbibliography[heading=none]
\end{frame}


\end{document}
%%% Local Variables:
%%% coding: utf-8
%%% mode: latex
%%% TeX-engine: luatex
%%% End:
